\chapter{合卷一：道路、册籍与可移动的帝国}
\label{ch:ming-synthesis-movement}

明初最醒目的景象是城墙、宫殿和北伐，但一个王朝能否维持，并不取决于都城是否宏伟。应天建国、北京营建、运河北运、卫所屯田、黄册鱼鳞册和边地设治，处理的是同一个难题：怎样让粮、兵、人和消息以可预期的方式移动。每一个装置都把距离缩短一点，也把新的成本交给沿线的人。

\section{把首都接在南方粮道上}

北方都城的选择并非只是一项军事姿态。宫城、仓储、坛庙、卫戍与漕运必须同时运作，才能让皇帝在北方持续发号。运河的疏浚和维护因而不是背景工程：闸坝、役夫、河患、仓场与江南赋粮把京师的日常供给接在长距离的征解链上。中央文书擅长记录额数、命令和河工方案，不能替代一段河道的淤塞、一户人家的劳役或一次运输的损耗。

\section{把人户接在责任链上}

黄册、鱼鳞图册、里甲和卫所试图把人口、土地、军役和赋役写进可追查的格式。册籍并不是社会本身；它是国家要求地方以某种方式说明社会的技术。书手、业主、军户、流民和地方官都可能在登记、隐匿、申报和催征中改变它的意义。正因如此，早明的“国家能力”既能看作建立比较和征收的能力，也必须看作不断同迁徙、土地转移和地方中介搏斗的能力。

\section{把边地接在补给与交换上}

从北元草原到云南、交趾、哈密与东南海面，明廷反复发现军事行动不能单独制造稳定边界。远征需要粮道，卫所需要军户，册封需要地方盟友，海禁也需要港口社会配合。交趾撤军不是简单的退却，而是承认直接统治的征粮、道路和抵抗成本；土木危机不是孤立的战场事故，而是北方警报、朝贡贸易、军饷与指挥失调集中爆裂的时刻。

\section{把非常时刻交还给日常秩序}

北京保卫、景泰即位和英宗复位显示，危机中的政治合法性离不开城门、军营、奏疏和储位等具体设施。一次非常授权可以救急，却必须随后决定谁掌兵、谁能定罪、谁的记忆被写进实录。早中明留下的制度得失正在这里：更密的道路、册籍和信息网络提高了调度能力，也让一处粮道、一份账册或一个宫门的失灵具有更远的后果。下一合卷将看到，白银和军费如何把这些连接推入新的规模。
